We now consider the case $\noisestd > 0$. The posterior density we consider in this section is given by
\[ 
    \noisyfilter{0}{\ltrmeas}{\ltrstate_0} \propto \pot{0}{\ltrmeas}{\ltrtopstate_0} \bwmarg{0}{\ltrstate_0}, \quad \mathrm{where} \quad \pot{0}\ltrmeas{\ltrtopstate_{0}} = \prod_{i = 1}^{\dimtr} \normpdf(\vecidx{\ltrmeas}{}{i}; \vecidx{\ltrtopstate}{0}{i}, (\noisestd / s_i)^2) \eqsp.
\]
Note that $\vecidx{\trtopstate}{\tau_i}{i} \sim \ndist{\alpha_{\tau_i}^{1/2} \vecidx{\ltrtopstate}{0}{i}}{1 - \alpha_{\tau_i}}$, and then $\normpdf(\alpha_{\tau_i}^{1/2}\vecidx{\ltrmeas}{}{i}; \alpha_{\tau_i}^{1/2} \vecidx{\ltrtopstate}{0}{i}, 1 - \alpha_{\tau_i})$ can be interpreted as the likelihood of observing $\alpha_{\tau_i}^{1/2}\vecidx{\ltrmeas}{}{i}$ for $\vecidx{\trtopstate}{\tau_i}{i}$.
This is generalization of \cref{eq:potential:coordinate:noisy} but where each $i$-th coordinate can have a possibly different $\tau_i$.
Let $\{ \tau_i \}_{i = 1} ^\dimtr \subset [1:n]$ with $\tau_1 < \dotsc < \tau_\dimtr$. For each $i \in [1:\top]$, define
\begin{equation}
    \label{eq:potential:coordinate:noisy}
    \pot{s, i}{\ltrmeas}{\vecidx{\ltrtopstate}{s}{i}} = \normpdf(\vecidx{\ltrtopstate}{s}{i}; \alpha_{s}^{1/2} \vecidx{\ltrmeas}{}{i}, (1 - (1 - \sigma^2 _\delta) \alpha_s / \alpha_{\tau_i})) 
\end{equation}
for $s \in [\tau_i, n]$, where $\sigma_{\delta}^2 > 0$ is a free hyper-parameter as in \cref{sec:noisy}. In the particular case of $\sigma^2 _\delta = 0$, $\pot{s,i}{\ltrmeas}{}$ is the density of forward process from $\tau_i$ to $s > \tau_i$ starting from $\alpha^{1/2} _{\tau_i} \vecidx{\ltrmeas}{}{i}$.
Then, define for $s \in [\tau_1:n]$, $\pot{s}{\ltrmeas}{\ltrtopstate_s} = \prod_{i = 1}^{\tau(s)}\pot{s, i}{\ltrmeas}{\vecidx{\ltrtopstate}{s}{i}}$ where $\tau(s) \eqdef \max \{ k \in [1:\dimtr]: s - \tau_k \geq 0\}$.

Denote for $\ell \in \intvect{1}{\dimx}$, $\projidx{\ltrstate}{}{\ell} \in \R^{\dimx-1}$ the vector $\ltrstate$ with its $\ell$-th coordinate removed.
The following proposition is a generalization of \cref{eq:noiseless_at_tau}.
\begin{proposition}
    \label{prop:noisytonoiseless}
    Assume that $\bwmarg{s+1}{\ltrstate_{s+1}} \bwtrans{}{s}{\ltrstate_{s+1}}{\ltrstate_s} = \bwmarg{s}{\ltrstate_s} \fwtrans{s+1}{\ltrstate_s}{\ltrstate_{s+1}}$ for all $s \in [0:n-1]$. Then it holds that
    \begin{equation*}
        \noisyfilter{0}{\ltrmeas}{\ltrstate_{0}} \propto
        \int \bwmarg{\tau_\dimtr}{\ltrstate_{\tau_\dimtr}} \delta_{\diffltrmeas_\dimtr}(\rmd \vecidx{\ltrstate}{\tau_\dimtr}{\dimtr}) \rmd \projidx{\ltrstate}{\tau_\dimtr}{\dimtr} \left\{ \prod_{i = 1}^{\dimtr-1} \bwtrans{}{\tau_i | \tau_{i + 1}}{\ltrstate_{\tau_{i+1}}}{\ltrstate_{\tau_{i}}} \delta_{\diffltrmeas_i}(\rmd \vecidx{\ltrstate}{\tau_i}{i}) \rmd \projidx{\ltrstate}{\tau_i}{i} \right\} \bwtrans{}{0|\tau_1}{\ltrstate_{\tau_1}}{\ltrstate_0} \eqsp.
    \end{equation*}
\end{proposition}
The proof of \Cref{prop:noisytonoiseless}  can be found in \Cref{proof:noisytonoiseless}. As in \cref{sec:noisy}, we can consider
the noisy problem as a sequence of noiseless observations of the variables $\trstate_{\tau_i}$ for $i \in \intvect{1}{\dimtr}$. Thus,
for $s \in \intvect{\tau_i}{\tau_{i+1}}$ we consider the potentials generated by all the observations $\lstate_{\tau_1}, \cdots, \lstate_{\tau_i}$.
Once the potentials are defined \eqref{eq:potential:coordinate:noisy}, we derive the expression for the proposal kernel $\bwopt{s}{\ltrmeas, \ell}{}{}$ and weight functions $\uweight{\ell}{s}$.
Consider $\ell \in \intvect{1}{\dimtr}$, $s \in \intvect{\tau_\ell}{n}$ and $\sigma_{\delta} > 0$. Define $\sigma^2 _{s|\tau_\ell} = 1 - (1 - \sigma_{\delta}^2)\alpha_{s} / \alpha_{\tau_\ell}$ and write $\overline{\boldsymbol{\mu}}_{s|s+1}(\ltrstate_{s+1}) \eqdef \boldsymbol{\mu}_{s+1}(\ltopstate_{s+1}, \topxpred{0|s+1}(\m{V} \ltrstate_{s+1}))$, $\underline{\boldsymbol{\mu}}_{s|s+1}(\ltrstate_{s+1}) \eqdef \boldsymbol{\mu}_{s+1}(\lbottomstate_{s+1}, \bottomxpred{0|s+1}(\m{V} \ltrstate_{s+1}))$ and so $\cat{\overline{\boldsymbol{\mu}}_{s|s+1}(\ltrstate_{s+1})}{\underline{\boldsymbol{\mu}}_{s|s+1}(\ltrstate_{s+1})}$ is the mean of the backward kernel $\bwtrans{}{s}{}{}$.
The density of the $\ell$-th coordinate of the proposal kernel is
\begin{equation*}
  \bwopt{s}{\ltrmeas, \ell}{\ltrstate_{s+1}}{\vecidx{\ltrtopstate}{s}{\ell}}
    \propto \normpdf\left(\alpha_{s}^{1/2}\vecidx{\ltrmeas}{}{\ell}; \vecidx{\ltrtopstate}{s}{\ell}, \sigma^2 _{s|\tau_\ell}\right)
    \normpdf\left(\vecidx{\ltrtopstate}{s}{\ell}; \overline{\boldsymbol{\mu}}_{s|s+1}(\ltrstate_{s+1})[\ell], \sigma_{s+1}^2  \right) \eqsp,
\end{equation*}
and the weight function is $\uweight{}{s}(\ltrstate_{s+1}) = \prod_{\ell=\tau(s)}^{1} \uweight{\ell}{s}(\ltrstate_{s+1})$ (we recall that $\tau(s) = \max \{ k \in [1:\dimtr]: s - \tau_k \geq 0\}$) where
\begin{equation*}
    \uweight{\ell}{s}(\ltrstate_{s+1}) \\
    \propto \frac{\int\normpdf\left(\alpha_{s}^{1/2}\vecidx{\ltrmeas}{}{\ell}; \vecidx{\ltrtopstate}{s}{\ell}, \sigma^2 _{s|\tau_\ell} \right) \normpdf\left(\vecidx{\ltrtopstate}{s}{\ell}; \overline{\boldsymbol{\mu}}_{s|s+1}(\ltrstate_{s+1})[\ell], \sigma_{s+1}^2  \right)}{\normpdf\left(\alpha_{s+1}^{1/2}\vecidx{\ltrmeas}{}{\ell}; \vecidx{\ltrtopstate}{s+1}{\ell}, \sigma^2 _{s|\tau_\ell}\right)} \eqsp.
\end{equation*}
Therefore, letting $\mathsf{k}_{s,\ell} \eqdef \sigma^2 _{s+1} / (\sigma^2 _{s+1} + \sigma^2 _{s|\tau_\ell})$ and $\tilde{\sigma}_{s, \ell}^2 \eqdef \sigma^2 _{s|\tau_\ell} \mathsf{k}_{s,\ell}$, we have
\begin{align*}
    \bwopt{s}{\ltrmeas, \ell}{\ltrstate_{s+1}}{\vecidx{\ltrtopstate}{s}{\ell}} & = \normpdf\left(\vecidx{\ltrtopstate}{s}{\ell}; \mathsf{k}_{s,\ell} \alpha^{1/2} _s \vecidx{\ltrmeas}{}{\ell} + (1 - \mathsf{k}_{s,\ell}) \overline{\boldsymbol{\mu}}_{s|s+1}(\ltrstate_{s+1})[\ell], \tilde{\sigma}_{s, \ell}^2\right) \eqsp,\\
    \uweight{\ell}{s}(\ltrstate_{s+1}) & = \frac{\normpdf\left(\alpha_{s}^{1/2}\vecidx{\ltrmeas}{}{\ell}; \overline{\boldsymbol{\mu}}_{s|s+1}(\ltrstate_{s+1})[\ell], \sigma_{s+1}^2 + \sigma^2 _{s|\tau_\ell}\right)}{\normpdf\left(\alpha_{s+1}^{1/2}\vecidx{\ltrmeas}{}{\ell}; \vecidx{\ltrtopstate}{s+1}{\ell}, \sigma^2 _{s+1|\tau_\ell}\right)} \eqsp.
\end{align*}
The proposal kernel is thus
\[
  \bwopt{s}{\ltrmeas}{\ltrstate_{s+1}}{\ltrstate_s} = \bwtrans{}{s}{\ltrstate_{s+1}}{\ltrbottomstate_s} \prod_{k = \tau(s) +1}^\dimtr \bwopt{s}{k}{\ltrstate_{s+1}}{\vecidx{\ltrstate}{s}{k}} \prod_{\ell=1}^{\tau(s)} \bwopt{s}{\ltrmeas, \ell}{\ltrstate_{s+1}}{\vecidx{\ltrtopstate}{s}{\ell}} \eqsp.
  \]
% \Cref{alg:algonoisygeneral} is an extension of \cref{alg:algonoisy} to this case.

% \begin{algorithm2e}[H]
%     \caption{\algo\ ($\noisestd > 0$)}
%     \label{alg:algonoisygeneral}
% \end{algorithm2e}
